\documentclass[10pt,a4paper]{amsart}
\usepackage{amsmath,amssymb,amsthm,graphicx}
\usepackage{verbatim}
\newenvironment{reference}{\comment}{\endcomment}
%\newenvironment{reference}{}{}
\newenvironment{slogan}{\comment}{\endcomment}
\newenvironment{history}{\comment}{\endcomment}
% For commutative diagrams we use Xy-pic
\usepackage[all]{xy}

% We use 2cell for 2-commutative diagrams.
\xyoption{2cell}
% Theorem environments.
%
\theoremstyle{plain}
\newtheorem{theorem}[subsection]{Theorem}
\newtheorem{proposition}[subsection]{Proposition}
\newtheorem{lemma}[subsection]{Lemma}

\theoremstyle{definition}
\newtheorem{definition}[subsection]{Definition}
\newtheorem{example}[subsection]{Example}
\newtheorem{exercise}[subsection]{Exercise}
\newtheorem{situation}[subsection]{Situation}

\theoremstyle{remark}
\newtheorem{remark}[subsection]{Remark}
\newtheorem{remarks}[subsection]{Remarks}

\numberwithin{equation}{subsection}

% Macros
%
\def\lim{\mathop{\mathrm{lim}}\nolimits}
\def\colim{\mathop{\mathrm{colim}}\nolimits}
\def\Spec{\mathop{\mathrm{Spec}}}
\def\Hom{\mathop{\mathrm{Hom}}\nolimits}
\def\Ext{\mathop{\mathrm{Ext}}\nolimits}
\def\SheafHom{\mathop{\mathcal{H}\!\mathit{om}}\nolimits}
\def\SheafExt{\mathop{\mathcal{E}\!\mathit{xt}}\nolimits}
\def\Sch{\mathit{Sch}}
\def\Mor{\mathop{\mathrm{Mor}}\nolimits}
\def\Ob{\mathop{\mathrm{Ob}}\nolimits}
\def\Sh{\mathop{\mathit{Sh}}\nolimits}
\def\NL{\mathop{N\!L}\nolimits}
\def\CH{\mathop{\mathrm{CH}}\nolimits}
\def\proetale{{pro\text{-}\acute{e}tale}}
\def\etale{{\acute{e}tale}}
\def\QCoh{\mathit{QCoh}}
\def\Ker{\mathop{\mathrm{Ker}}}
\def\Im{\mathop{\mathrm{Im}}}
\def\Coker{\mathop{\mathrm{Coker}}}
\def\Coim{\mathop{\mathrm{Coim}}}

% Boxtimes
%
\DeclareMathSymbol{\boxtimes}{\mathbin}{AMSa}{"02}

%
% Macros for moduli stacks/spaces
%
\def\QCohstack{\mathcal{QC}\!\mathit{oh}}
\def\Cohstack{\mathcal{C}\!\mathit{oh}}
\def\Spacesstack{\mathcal{S}\!\mathit{paces}}
\def\Quotfunctor{\mathrm{Quot}}
\def\Hilbfunctor{\mathrm{Hilb}}
\def\Curvesstack{\mathcal{C}\!\mathit{urves}}
\def\Polarizedstack{\mathcal{P}\!\mathit{olarized}}
\def\Complexesstack{\mathcal{C}\!\mathit{omplexes}}
% \Pic is the operator that assigns to X its picard group, usage \Pic(X)
% \Picardstack_{X/B} denotes the Picard stack of X over B
% \Picardfunctor_{X/B} denotes the Picard functor of X over B
\def\Pic{\mathop{\mathrm{Pic}}\nolimits}
\def\Picardstack{\mathcal{P}\!\mathit{ic}}
\def\Picardfunctor{\mathrm{Pic}}
\def\Deformationcategory{\mathcal{D}\!\mathit{ef}}
\usepackage{cleveref}
\crefname{originallemma}{original lemma}{original lemmas}
\Crefname{originallemma}{Original lemma}{Original lemmas}
\makeatletter
\newlabel{lemma-coherent-Noetherian}{{07UB}{source}{}{sourcepointer-1}{}}
\newlabel{lemma-coherent-Noetherian@cref}{{[originallemma][0][]07UB}{[1][0][]source}{}{}{}}
\makeatother
\crefname{originallemma}{original lemma}{original lemmas}
\Crefname{originallemma}{Original lemma}{Original lemmas}
\makeatletter
\newlabel{lemma-directed-colimit-coherent}{{07UV}{source}{}{sourcepointer-2}{}}
\newlabel{lemma-directed-colimit-coherent@cref}{{[originallemma][0][]07UV}{[1][0][]source}{}{}{}}
\makeatother
\crefname{originallemma}{original lemma}{original lemmas}
\Crefname{originallemma}{Original lemma}{Original lemmas}
\makeatletter
\newlabel{lemma-finite-presentation-quasi-compact-colimit}{{07U7}{source}{}{sourcepointer-3}{}}
\newlabel{lemma-finite-presentation-quasi-compact-colimit@cref}{{[originallemma][0][]07U7}{[1][0][]source}{}{}{}}
\makeatother
\crefname{originallemma}{original lemma}{original lemmas}
\Crefname{originallemma}{Original lemma}{Original lemmas}
\makeatletter
\newlabel{spaces-properties-lemma-pushforward-etale-base-change-modules}{{03LX}{source}{}{sourcepointer-4}{}}
\newlabel{spaces-properties-lemma-pushforward-etale-base-change-modules@cref}{{[originallemma][0][]03LX}{[1][0][]source}{}{}{}}
\makeatother
\crefname{originalsection}{original section}{original sections}
\Crefname{originalsection}{Original section}{Original sections}
\makeatletter
\newlabel{spaces-properties-section-properties-modules}{{05VR}{source}{}{sourcepointer-5}{}}
\newlabel{spaces-properties-section-properties-modules@cref}{{[originalsection][0][]05VR}{[1][0][]source}{}{}{}}
\makeatother
\title{On an affine space over a Noetherian base every quasi-coherent module is a filtered colimit of finitely presented modules}
\author{The Stacks Project Authors}
\date{stacks-project:spaces-cohomology.tex:local-snapshot-20261012; GNU FDL 1.2 or later}
\begin{document}
\maketitle
\noindent\textit{Editorial note.} All delivered mathematical text is the author's own, copied from the pinned Stacks Project chapter file spaces-cohomology.tex; no statement or proof step was rewritten, reordered or removed.\par\medskip
\noindent\textit{Editorial note.} Preamble provenance: the pinned chapter file opens with \textbackslash{}input\{preamble\} and carries no preamble text of its own, so the project's single Stacks-wide preamble staged at sources/stacks-spaces-cohomology/preamble.tex is excerpted verbatim as three spans, lines 8-12 (verbatim package and the project's own reference/slogan/history display environments), lines 14-18 (Xy-pic setup) and lines 154-217 (theorem declarations and mathematical macros). The wrapper class and the amsmath/amssymb/amsthm/graphicx base packages are supplied by the assembler. No macro was invented or redefined.\par\medskip
\noindent\textit{Editorial note.} Same-chapter pointers are located at their \textbackslash{}label lines inside the pinned chapter file; cross-chapter pointers are located in the project's own tag list staged in the snapshot (sources/stacks-chow/tags.txt), at the verbatim lines carrying their canonical tags.\par\medskip
\noindent\textit{Allowed standard prerequisites (editorial).}\begin{itemize}
\item Every quasi-coherent module on a Noetherian space is a filtered colimit of coherent modules. Same chapter spaces-cohomology.tex, label lemma-directed-colimit-coherent (invoked at proof line 2909). It writes \$f\_*\textbackslash{}mathcal\{F\} = \textbackslash{}colim \textbackslash{}mathcal\{H\}\_i\$ with \$\textbackslash{}mathcal\{H\}\_i\$ coherent, which is the system of finitely presented modules that is pulled back.
\item Coherence on a Noetherian affine base is finite presentation. Same chapter spaces-cohomology.tex, label lemma-coherent-Noetherian (invoked at proof line 2910). It upgrades the coherent modules \$\textbackslash{}mathcal\{H\}\_i\$ to modules of finite presentation, so that their pullbacks are finitely presented as well.
\item Base change of pushforwards and pullbacks along an etale cover, which reduces the surjectivity claim to schemes. Properties of Spaces, label spaces-properties-lemma-pushforward-etale-base-change-modules (invoked at proof line 2924). With \$U = X \textbackslash{}times\_Y V\$ it identifies the restriction of \$f\textasciicircum{}*f\_*\textbackslash{}mathcal\{F\} \textbackslash{}to \textbackslash{}mathcal\{F\}\$ with the corresponding map for the affine morphism of schemes \$f' : U \textbackslash{}to V\$, where surjectivity is known.
\item A morphism between filtered colimits of finitely presented modules factors through a stage. Same chapter spaces-cohomology.tex, label lemma-finite-presentation-quasi-compact-colimit (invoked at proof line 2953). For every \$j\$ it produces a stage \$i\$ and a morphism \$\textbackslash{}alpha : \textbackslash{}mathcal\{G\}\_j \textbackslash{}to \textbackslash{}mathcal\{H\}\_i\$ making the obvious square commute, which is what builds the directed set of triples.
\end{itemize}\medskip
\ifdefined\hypertarget\hypertarget{sourcepointer-1}{}\fi
\noindent\textit{Source reference.} Original lemma 07UB: sources/stacks-spaces-cohomology/pinned-original-source.tex, source lines 2094--2094. Navigation pointer for \textbackslash{}ref\{lemma-coherent-Noetherian\} used inside the counted span (line 2910); the argument that consumes it is retained in the span.\par\medskip
\ifdefined\hypertarget\hypertarget{sourcepointer-2}{}\fi
\noindent\textit{Source reference.} Original lemma 07UV: sources/stacks-spaces-cohomology/pinned-original-source.tex, source lines 2866--2866. Navigation pointer for \textbackslash{}ref\{lemma-directed-colimit-coherent\} used inside the counted span (line 2909); the argument that consumes it is retained in the span.\par\medskip
\ifdefined\hypertarget\hypertarget{sourcepointer-3}{}\fi
\noindent\textit{Source reference.} Original lemma 07U7: sources/stacks-spaces-cohomology/pinned-original-source.tex, source lines 735--735. Navigation pointer for \textbackslash{}ref\{lemma-finite-presentation-quasi-compact-colimit\} used inside the counted span (line 2953); the argument that consumes it is retained in the span.\par\medskip
\ifdefined\hypertarget\hypertarget{sourcepointer-4}{}\fi
\noindent\textit{Source reference.} Original lemma 03LX: sources/stacks-chow/tags.txt, source lines 4729--4729. Navigation pointer for \textbackslash{}ref\{spaces-properties-lemma-pushforward-etale-base-change-modules\} used inside the counted span (line 2924); the argument that consumes it is retained in the span.\par\medskip
\ifdefined\hypertarget\hypertarget{sourcepointer-5}{}\fi
\noindent\textit{Source reference.} Original section 05VR: sources/stacks-chow/tags.txt, source lines 7560--7560. Navigation pointer for \textbackslash{}ref\{spaces-properties-section-properties-modules\} used inside the counted span (line 2914); the argument that consumes it is retained in the span.\par\medskip
\begin{lemma}
\label{lemma-direct-colimit-finite-presentation}
Let $S$ be a scheme. Let $f : X \to Y$ be an affine morphism of algebraic
spaces over $S$ with $Y$ Noetherian. Then every quasi-coherent
$\mathcal{O}_X$-module is a filtered colimit of finitely presented
$\mathcal{O}_X$-modules.
\end{lemma}

\begin{proof}
Let $\mathcal{F}$ be a quasi-coherent $\mathcal{O}_X$-module.
Write $f_*\mathcal{F} = \colim \mathcal{H}_i$ with $\mathcal{H}_i$
a coherent $\mathcal{O}_Y$-module, see
Lemma \ref{lemma-directed-colimit-coherent}.
By Lemma \ref{lemma-coherent-Noetherian} the modules $\mathcal{H}_i$
are $\mathcal{O}_Y$-modules of finite presentation. Hence
$f^*\mathcal{H}_i$ is an $\mathcal{O}_X$-module of finite presentation, see
Properties of Spaces, Section
\ref{spaces-properties-section-properties-modules}.
We claim the map
$$
\colim f^*\mathcal{H}_i = f^*f_*\mathcal{F} \to \mathcal{F}
$$
is surjective as $f$ is assumed affine, Namely, choose
a scheme $V$ and a surjective \'etale morphism $V \to Y$. Set
$U = X \times_Y V$. Then $U$ is a scheme, $f' : U \to V$ is affine, and
$U \to X$ is surjective \'etale. By
Properties of Spaces, Lemma
\ref{spaces-properties-lemma-pushforward-etale-base-change-modules}
we see that $f'_*(\mathcal{F}|_U) = f_*\mathcal{F}|_V$ and similarly
for pullbacks. Thus the restriction of $f^*f_*\mathcal{F} \to \mathcal{F}$
to $U$ is the map
$$
f^*f_*\mathcal{F}|_U = (f')^*(f_*\mathcal{F})|_V) =
(f')^*f'_*(\mathcal{F}|_U) \to \mathcal{F}|_U
$$
which is surjective as $f'$ is an affine morphism of schemes.
Hence the claim holds.

\medskip\noindent
We conclude that every quasi-coherent module on $X$ is a quotient of a
filtered colimit of finitely presented modules. In particular, we see that
$\mathcal{F}$ is a cokernel of a map
$$
\colim_{j \in J} \mathcal{G}_j \longrightarrow \colim_{i \in I} \mathcal{H}_i
$$
with $\mathcal{G}_j$ and $\mathcal{H}_i$ finitely presented. Note
that for every $j \in I$ there exist $i \in I$ and a morphism
$\alpha : \mathcal{G}_j \to \mathcal{H}_i$ such that
$$
\xymatrix{
\mathcal{G}_j \ar[r]_\alpha \ar[d] & \mathcal{H}_i \ar[d] \\
\colim_{j \in J} \mathcal{G}_j \ar[r] &
\colim_{i \in I} \mathcal{H}_i
}
$$
commutes, see
Lemma \ref{lemma-finite-presentation-quasi-compact-colimit}.
In this situation $\Coker(\alpha)$ is a finitely presented
$\mathcal{O}_X$-module which comes endowed with a map
$\Coker(\alpha) \to \mathcal{F}$. Consider the set $K$ of
triples $(i, j, \alpha)$ as above. We say that
$(i, j, \alpha) \leq (i', j', \alpha')$ if and only if
$i \leq i'$, $j \leq j'$, and the diagram
$$
\xymatrix{
\mathcal{G}_j \ar[r]_\alpha \ar[d] & \mathcal{H}_i \ar[d] \\
\mathcal{G}_{j'} \ar[r]^{\alpha'} &
\mathcal{H}_{i'}
}
$$
commutes. It follows from the above that $K$ is
a directed partially ordered set,
$$
\mathcal{F} = \colim_{(i, j, \alpha) \in K} \Coker(\alpha),
$$
and we win.
\end{proof}

\end{document}
